\uuid{VmLM}
\titre{Matrice/Base/Dimension}
\auteur{hanani ; mzari}
\organisation{exo7}
\licence{CC BY-NC-SA 4.0}
\datecreate{2020-01-20}
\niveau{L1}
\module{Algèbre}
\chapitre{Matrice}
\sousChapitre{Propriétés élémentaires, généralités}
\difficulte{1}
\theme{Calcul matriciel, Matrice, Base, Dimension}
\exo7qcmid{lille-2.matrices.004}

\contenu{
\begin{question}
\qclassification{108.01}
\qtags{motcle=Matrice/Base/Dimension}
On considère $M_{n,m} (\Rr)$ l'ensemble des matrices à $n$ lignes et $m$ colonnes, à  coefficients dans $\Rr$, muni de l'addition usuelle et la multiplication par un scalaire. Quelles sont les assertions vraies ?
\begin{answers}
\good{$M_{n,m} (\Rr)$ est un espace vectoriel}.
\good{$\dim M_{n,m} (\Rr) = mn $}.
\bad{$\dim M_{n,m} (\Rr) = m+n $}.
\bad{$M_{n,m} (\Rr)$ est un espace vectoriel de dimension infinie}.
\end{answers}
\begin{explanations} On vérifie que $M_{n,m} (\Rr)$, muni des opérations usuelles est un $\Rr$- espace vectoriel.
\vskip0mm
Pour $1\le i\le n$ et $1\le j\le m$, on note $D_{i,j}$ la matrice dont le coefficient située à la ième ligne et jième colonne est $1$ et les autres coefficients sont nuls. Alors, $\{D_{i,j}\; ; \; 1\le i\le n, 1\le j\le m\}$ est une base de $M_{n,m} (\Rr)$. Par conséquent, $\dim M_{n,m} (\Rr) = mn $.
\end{explanations}
\end{question}
}
